\chapter{Introdução}
\label{cap:introducao}

\section{Contextualização e problema}
\label{sec:contextualizacao-problema}

Aplicações móveis em saúde, frequentemente denominadas mHealth, têm sido discutidas como ferramentas de apoio ao acesso a informações de saúde, ao acompanhamento de pacientes e à organização de atividades relacionadas ao cuidado. Por meio de dispositivos móveis, essas aplicações podem ampliar a disponibilidade de informações, apoiar a comunicação de orientações e favorecer a continuidade do cuidado fora do ambiente clínico. Nesse contexto, sistemas mHealth podem complementar o acompanhamento realizado por profissionais de saúde, especialmente em cenários que envolvem tratamentos contínuos ou recorrentes e planos de cuidado estruturados \cite{who2011mhealth,emerson2022mhealth,niemiec2022takere}.

Apesar desse potencial, a efetividade de aplicações mHealth depende não apenas da disponibilidade de informações, mas também da capacidade dos usuários de compreendê-las e aplicá-las em seu cotidiano. Muitos conteúdos de saúde são tradicionalmente produzidos com linguagem técnica, uso de jargões médicos, frases extensas e estruturas textuais pouco adequadas ao público leigo. Quando essas informações são transferidas para interfaces digitais sem adaptação comunicacional, existe o risco de que o paciente tenha dificuldade para interpretar orientações, compreender alertas ou reconhecer ações esperadas dentro de seu plano de cuidado.

Esse problema está diretamente relacionado ao conceito de letramento em saúde, que envolve a capacidade de acessar, compreender, avaliar e utilizar informações de saúde para tomar decisões adequadas no dia a dia. Em aplicações mHealth, essa questão torna-se ainda mais relevante, pois usuários com diferentes níveis de escolaridade, familiaridade tecnológica e conhecimento sobre saúde podem interagir com os mesmos conteúdos. Emerson et al. \cite{emerson2022mhealth} destacam que a entrega de informações de saúde por dispositivos móveis amplia o acesso a conteúdos de saúde, mas também cria desafios para garantir que esses materiais sejam adequados a diferentes necessidades de letramento.

A comunicação em linguagem simples é uma estratégia importante para enfrentar esse desafio. Materiais de saúde escritos de forma clara, objetiva e orientada ao público-alvo tendem a facilitar a compreensão e o uso das informações apresentadas. Assim, a simplificação textual em saúde não consiste apenas em reduzir a complexidade gramatical, mas em tornar a informação mais acessível, buscando preservar seu significado, sua precisão e sua utilidade prática \cite{cdc2025plainlanguage}.

Este trabalho será contextualizado no ambiente da Takere, uma plataforma no-code para criação de aplicações mHealth baseadas em planos de cuidado. A Takere permite que profissionais de saúde instanciem aplicações móveis para pacientes a partir de informações de seus planos de cuidado, reduzindo a dependência de equipes de desenvolvimento para a criação de soluções específicas \cite{niemiec2022takere}. Neste TCC, a Takere não constitui o tema central da pesquisa, mas funciona como ambiente de referência para investigar como conteúdos técnicos presentes em planos de cuidado podem ser simplificados e adaptados linguisticamente para pacientes.

\section{Justificativa}
\label{sec:justificativa}

A relevância deste trabalho decorre do papel que a compreensão das informações de saúde exerce no uso adequado de orientações de cuidado. Em aplicações mHealth baseadas em planos de cuidado, textos, alertas e explicações orientam ações que o paciente deve interpretar e realizar fora do ambiente clínico. Assim, quando esses conteúdos são apresentados de forma excessivamente técnica ou pouco adequada ao perfil do usuário, a comunicação pode se tornar uma barreira para o acompanhamento do cuidado.

Esse ponto é especialmente importante porque a compreensão das informações de saúde pode influenciar a capacidade do paciente de reconhecer instruções, interpretar alertas e seguir orientações no cotidiano. Estudos sobre letramento em saúde indicam associação positiva entre letramento em saúde e adesão ao tratamento, embora essa relação seja multifatorial. Miller \cite{miller2016healthliteracy} identificou associação positiva entre letramento em saúde e adesão ao tratamento, enquanto Hyvert et al. \cite{hyvert2023healthliteracy} destacam que o letramento em saúde tem papel relevante na adesão medicamentosa em doenças crônicas, embora fatores como crenças, autoeficácia, conhecimento sobre medicamentos e condições sociodemográficas também devam ser considerados.

Nesse sentido, o trabalho justifica-se por investigar uma solução funcional voltada à adaptação linguística de conteúdos técnicos de saúde, sem modificar condutas clínicas, gerar diagnósticos ou criar novas recomendações. A proposta posiciona os LLMs como apoio à comunicação e à compreensão do paciente, exigindo o uso de guardrails para preservar a fidelidade das informações originais e evitar alterações em conteúdos clinicamente relevantes.

\section{Proposta do trabalho}
\label{sec:proposta-trabalho}

Diante desse cenário, este trabalho propõe o uso de Modelos de Linguagem de Grande Porte, ou Large Language Models (LLMs), para simplificar e adaptar textos técnicos de saúde apresentados em plataformas mHealth. A solução funcional desenvolvida no contexto do TCC recebe textos ou arquivos contendo orientações de saúde previamente definidas ou validadas por profissionais e gera versões mais claras, considerando características do perfil do paciente.

Esse perfil pode incluir idade, escolaridade, nível de letramento em saúde e condição médica geral, com o objetivo de ajustar vocabulário, tom, nível de detalhamento e forma de explicação. A adaptação proposta é estritamente comunicacional e linguística, não clínica. Portanto, a solução não deve gerar diagnósticos, prescrever tratamentos, criar novas recomendações médicas ou substituir a atuação de profissionais de saúde.

O uso de LLMs nesse contexto é motivado pela capacidade desses modelos de processar e gerar linguagem natural, o que permite reescrever conteúdos, explicar termos técnicos e adaptar materiais educacionais a diferentes públicos. Aydin et al. \cite{aydin2024llms}, em uma revisão de escopo sobre LLMs na educação de pacientes, identificam aplicações como geração de materiais educativos, interpretação de informações médicas sob a perspectiva do paciente, tradução de conteúdos e melhoria da comunicação entre pacientes e profissionais. De forma mais específica, Loughran et al. \cite{loughran2024mhealth} discutem o uso de LLMs para enfrentar problemas de letramento em saúde em mHealth, avaliando seu potencial para reduzir a complexidade linguística de materiais educacionais em saúde.

Entretanto, o uso de LLMs em saúde exige cautela. Embora esses modelos possam apoiar a simplificação de textos e a produção de materiais mais acessíveis, eles também podem gerar respostas imprecisas, pouco confiáveis ou inadequadas ao contexto de uso. Em conteúdos de saúde, tais falhas podem comprometer a qualidade da comunicação com o paciente, especialmente quando o usuário não possui conhecimento técnico suficiente para avaliar criticamente a informação apresentada. Aydin et al. \cite{aydin2024llms} destacam que, apesar do potencial dos LLMs na educação de pacientes, ainda existem desafios relacionados à acurácia, à legibilidade e à confiabilidade dos conteúdos gerados.

Por isso, a simplificação proposta neste trabalho deve buscar preservar o significado médico original, mantendo alertas, riscos, contraindicações, dosagens, frequências e instruções de segurança sem alterações indevidas. Nesse sentido, guardrails são tratados como mecanismos de controle aplicados ao processo de geração textual, com o objetivo de orientar e restringir a saída da LLM em relação ao conteúdo original. De forma geral, Dong et al. \cite{dong2024guardrails} descrevem guardrails como mecanismos voltados a filtrar entradas e saídas de LLMs, buscando reduzir riscos associados ao uso desses modelos. No contexto deste trabalho, esses mecanismos são utilizados com uma finalidade delimitada: apoiar a fidelidade da simplificação textual, evitando a criação de novas informações clínicas ou a modificação de instruções relevantes.

\section{Objetivo e pergunta de pesquisa}
\label{sec:objetivo-pergunta-pesquisa}

O objetivo geral deste trabalho é investigar, implementar e avaliar uma solução baseada em LLMs para simplificação e adaptação de linguagem técnica em uma plataforma mHealth baseada em planos de cuidado, utilizando guardrails para preservar a precisão médica das informações apresentadas ao paciente. A avaliação da solução considera tanto a preservação das informações médicas essenciais quanto indicadores de simplificação linguística, incluindo métricas automáticas de legibilidade e complexidade textual. A pesquisa busca, portanto, explorar o potencial de LLMs como apoio à comunicação em saúde digital, sem atribuir a esses modelos funções clínicas autônomas ou decisórias.

A partir desse objetivo, a pergunta de pesquisa que orienta o trabalho é:

\begin{quote}
Como Modelos de Linguagem de Grande Porte podem ser utilizados, com apoio de guardrails, para simplificar e adaptar textos técnicos de saúde ao perfil de pacientes, preservando a precisão médica das informações apresentadas?
\end{quote}

\section{Organização do trabalho}
\label{sec:organizacao-trabalho}

O restante do trabalho está organizado de forma a aprofundar os conceitos necessários para essa investigação, apresentar a solução desenvolvida, descrever os mecanismos utilizados para controle da geração textual e discutir sua avaliação. Inicialmente, são abordados os fundamentos relacionados a mHealth, letramento em saúde, linguagem simples, LLMs, guardrails e avaliação da legibilidade textual. Em seguida, é descrita a implementação funcional no contexto da plataforma-alvo, bem como os critérios adotados para preservar a fidelidade ao conteúdo original. Por fim, são apresentados os resultados, limitações e possibilidades de continuidade da pesquisa.
